% UNCHECKED HAND-PROOF SOURCE, 2026-11-15.
% No compilation, proof-assistant checking, or numerical verification was run.
% Input fragment using standard amsmath/amssymb notation.
\section{Complete inputs for the compatible-idempotent gap certificate}
\label{urgent:oblique}

\paragraph{Status and conventions.}
This expanded hand proof is prepared for formalization and has not been
machine checked. It derives the hypothesis called hcos rather than assuming
it. Hilbert spaces are finite dimensional over $\mathbb C$; norms of maps
are operator norms; $[A,B]=AB-BA$; $\Pi_M$ denotes the orthogonal projector
onto $M$. Define the Friedrichs cosine, including degenerate cases, by
\[
 c(M,N)=\|\Pi_M\Pi_N-\Pi_{M\cap N}\|.
\]
Writing $W=M\cap N$, the projector $\Pi_W$ commutes with both projectors.
Thus the operator inside this norm is
$(\Pi_M-\Pi_W)(\Pi_N-\Pi_W)$, and its norm equals the supremum of
$|\langle u,v\rangle|$ over unit vectors in $M\ominus W$ and $N\ominus W$.
The operator-norm definition gives zero when a reduced subspace is zero.

\subsection{The orthogonalizer and its commutators}

\paragraph{Lemma O1 (idempotent normal form).}
If $P^2=P$ and $M=\operatorname{ran}P$, the orthogonal splitting
$\mathcal H=M\oplus M^\perp$ gives
\[
 P=\begin{pmatrix}I&A\\0&0\end{pmatrix},\quad
 P^*=\begin{pmatrix}I&0\\A^*&0\end{pmatrix},\quad
 S:=P-P^*=\begin{pmatrix}0&A\\-A^*&0\end{pmatrix}.
\]
Indeed $P$ fixes its range, and all its outputs lie in its range.
Multiplication gives $S^*S=\operatorname{diag}(AA^*,A^*A)$, hence
$\|S\|=\|A\|$. For $u\in M$ and $v\in M^\perp$,
\[
 \|P(u,v)\|\leq\|u\|+\|A\|\|v\|
 \leq\sqrt{1+\|A\|^2}\sqrt{\|u\|^2+\|v\|^2}.
\]
Consequently, with $a=\|P-P^*\|$ and $\kappa=\sqrt{1+a^2}$,
$\|P\|\leq\kappa$. This inequality also covers $P=0$.

\paragraph{Lemma O2 (exact orthogonalizer).}
Set $D=I-(P-P^*)^2$. Since $S^*=-S$, $D=I+S^*S\geq I$.
Thus $D$ is invertible and $\|D^{-1}\|\leq1$. In the above blocks,
\[
 D=\operatorname{diag}(I+AA^*,I+A^*A),\qquad
 PP^*=\operatorname{diag}(I+AA^*,0).
\]
Therefore
\begin{equation}
 \Pi_M=PP^*D^{-1}.
 \label{urgent:orthogonalizer}
\end{equation}
The block computation identifies the right side with the actual orthogonal
projection; it does not assume that an arbitrary product is self-adjoint.

\paragraph{Lemma O3 (orthogonalization commutator estimate).}
For any map $T$ put $u_T=\|[P,T]\|+\|[P^*,T]\|$ and
$L=\kappa+2\kappa^2a$. Then
\begin{equation}
 \|[\Pi_M,T]\|\leq L u_T.
 \label{urgent:ortho-commutator}
\end{equation}
The algebraic identities are
\[
 [XY,T]=X[Y,T]+[X,T]Y,\qquad
 [D^{-1},T]=-D^{-1}[D,T]D^{-1}.
\]
The second follows by expanding $[DD^{-1},T]=0$.
For $N=PP^*$, the first gives
$\|[N,T]\|\leq\kappa u_T$ and $\|N\|\leq\kappa^2$.
Also $[S,T]=[P,T]-[P^*,T]$, so
\[
 \|[D,T]\|=\|[S^2,T]\|\leq2a u_T.
\]
Apply the product identity to $ND^{-1}$ and use $\|D^{-1}\|\leq1$:
the two terms are bounded by $\kappa u_T$ and $2\kappa^2a u_T$.

\paragraph{Lemma O4 (two compatible idempotents).}
Suppose $P^2=P$, $R^2=R$, and $PR=RP$. Write
$\Pi=\Pi_{\operatorname{ran}P}$, $\Theta=\Pi_{\operatorname{ran}R}$,
and $\epsilon=\|[P,R^*]\|$. Then
\begin{equation}
 \|[\Pi,\Theta]\|\leq2L_PL_R\epsilon.
 \label{urgent:double-commutator}
\end{equation}
Indeed $[P^*,R^*]=0$ and $[P,R^*]^*=[R,P^*]$.
Apply O3 to $R$ first with $T=P$, then with $T=P^*$.
It gives $\|[\Theta,P]\|,\|[\Theta,P^*]\|\leq L_R\epsilon$.
Applying O3 to $P$ with $T=\Theta$ now proves the claim.
The adjoint identity also proves the symmetry
$\|[P,R^*]\|=\|[R,P^*]\|$.

\subsection{Compatibility supplies angular separation}

\paragraph{Lemma O5 (quotient separation).}
Put $M=\operatorname{ran}P$, $N=\operatorname{ran}R$, $W=M\cap N$.
The product $PR$ has range $W$: it takes values in both ranges by
commutativity, and fixes every vector in their intersection.
Both maps preserve $W$. Their induced maps $\overline P,\overline R$
on $\mathcal H/W$ are idempotents, have product zero in either order, and
satisfy $\|\overline P\|\leq\|P\|$ and $\|\overline R\|\leq\|R\|$.
For the norm statement, for each representative $x$ and every $w\in W$,
\[
 \|[Px]\|_{\mathcal H/W}
 =\|[P(x+w)]\|_{\mathcal H/W}\leq\|P\|\|x+w\|;
\]
take the infimum over $w$. Identify the quotient isometrically with
$W^\perp$ using orthogonal representatives. The range of $\overline P$
then equals $M\ominus W$: for $m\in M$, its representative
$m-\Pi_Wm$ remains in $M$, and every image class has such a representative.
The corresponding assertion holds for $R$.

For nonzero $u$ and $v$ in these two quotient ranges,
$\overline P u=u$ and $\overline P v=0$, so for every $t\in\mathbb C$,
\[
 \|u\|\leq\kappa_P\|u-tv\|.
\]
Minimize the squared norm over $t$ to obtain
\[
 \|u\|^2\leq\kappa_P^2
 \left(\|u\|^2-\frac{|\langle u,v\rangle|^2}{\|v\|^2}\right).
\]
Thus
\begin{equation}
 c(M,N)\leq \frac{a_P}{\kappa_P},\qquad
 c(M,N)\leq \frac{a_R}{\kappa_R}.
 \label{urgent:quotient-separation}
\end{equation}
If either reduced range is zero the assertion follows from $c=0$.
Since $\kappa_P^2=1+a_P^2$, these estimates imply
\[
 \sqrt{1-c(M,N)^2}\geq
 \frac1{\min(\kappa_P,\kappa_R)}.
\]
In particular the positive denominator needed below is established;
small commutators of arbitrary orthogonal projections alone would not
give this exclusion of nearly parallel ranges.

\paragraph{Lemma O6 (cosine from an orthogonal commutator).}
For orthogonal projections $A,B$, put
$R=\Pi_{\operatorname{ran}A\cap\operatorname{ran}B}$ and
$c=\|AB-R\|$. Then
\begin{equation}
 c\sqrt{1-c^2}\leq\|[A,B]\|.
 \label{urgent:cosine-commutator}
\end{equation}
To see this without assuming a principal-angle theorem, restrict to
$W^\perp$ and put $A_0=A-R$, $B_0=B-R$.
The positive map $T=A_0B_0A_0$ has norm $c^2$, by
$\|A_0B_0\|^2=\|(A_0B_0)(A_0B_0)^*\|$.
If $c=0$ the claim is immediate. Otherwise let $u$ be a unit eigenvector
for $T$ with eigenvalue $c^2$. Since this eigenvalue is positive,
$A_0u=u$. Therefore
\[
 [A_0,B_0]u=c^2u-B_0u,\qquad
 \|B_0u\|^2=\langle u,B_0u\rangle=c^2.
\]
Expanding the squared norm gives $c^2-c^4$, which proves the bound.
Combine this with O4--O5 to conclude
\begin{equation}
 c(M,N)\leq
 2\min(\kappa_P,\kappa_R)L_PL_R\|[P,R^*]\|.
 \label{urgent:adjoint-cosine}
\end{equation}

\subsection{Derivation of the physical cross-term inequality}

\paragraph{Lemma O7 (complementary Friedrichs cosines).}
For finite-dimensional subspaces,
$c(M,N)=c(M^\perp,N^\perp)$.
Here are the details needed if this identity is not already available
in a proof-assistant library. For $A=\Pi_M$, $B=\Pi_N$, diagonalize
the positive contraction $ABA$ on $\operatorname{ran}A$.
An eigenvector of eigenvalue $1$ belongs to $M\cap N$.
An eigenvector of eigenvalue $0$ belongs to $M\cap N^\perp$, since
$\|Bu\|^2=\langle u,ABA u\rangle$.
For each unit eigenvector $u$ of eigenvalue $0<\lambda<1$, set
\[
 v=\frac{Bu-\lambda u}{\sqrt{\lambda(1-\lambda)}}.
\]
Then $Av=0$, $\|v\|=1$, and on the span of $u,v$,
\[
 A=\begin{pmatrix}1&0\\0&0\end{pmatrix},\qquad
 B=\begin{pmatrix}\lambda&\sqrt{\lambda(1-\lambda)}\\
 \sqrt{\lambda(1-\lambda)}&1-\lambda\end{pmatrix}.
\]
These formulas follow by $ABA u=\lambda u$ and $B^2=B$.
Different orthogonal eigenvectors produce orthogonal two-dimensional
blocks: use
$\langle Bu,Bu'\rangle=\langle u,Bu'\rangle
=\lambda'\langle u,u'\rangle$
and expand the numerators defining $v,v'$.
The same computation works within a repeated eigenspace.
Each block reduces both self-adjoint projections.
On the remaining reducing space $A=0$, so diagonalizing the remaining
projection $B$ gives only scalar blocks $(A,B)=(0,0),(0,1)$.
Together with the eigenvalue-zero and eigenvalue-one blocks already
listed, this exhausts the space.

On each nontrivial block the range cosine is $\sqrt\lambda$.
The complementary ranges are generated by $v$ and
$-\sqrt{1-\lambda}\,u+\sqrt\lambda\,v$, also with overlap
$\sqrt\lambda$. All scalar blocks contribute zero once the relevant
intersection is removed. The maximum over the blocks, or zero when
there are no such blocks, proves the claimed equality.

\paragraph{Lemma O8 (hcos is a conclusion).}
For $Q=I-\Pi_M$, $T=I-\Pi_N$ and every $\psi$,
\begin{equation}
 \operatorname{Re}\langle Q\psi,T\psi\rangle
 \geq-c(M,N)\|Q\psi\|\|T\psi\|.
 \label{urgent:hcos}
\end{equation}
Put $Z=M^\perp\cap N^\perp$ and $z=\Pi_Z\psi$.
The projections of both $Q\psi$ and $T\psi$ onto $Z$ equal $z$.
Thus write $Q\psi=z+u$, $T\psi=z+v$, with
$u\in M^\perp\ominus Z$ and $v\in N^\perp\ominus Z$.
O7 and the definition of the cosine give
$|\langle u,v\rangle|\leq c(M,N)\|u\|\|v\|$.
Consequently
\[
 \operatorname{Re}\langle Q\psi,T\psi\rangle
 =\|z\|^2+\operatorname{Re}\langle u,v\rangle
 \geq-c(M,N)\|u\|\|v\|
 \geq-c(M,N)\|Q\psi\|\|T\psi\|.
\]
Equivalently, the anticommutator quadratic form is at least
$-2c(M,N)\|Q\psi\|\|T\psi\|$.

\subsection{Construction of the matrix and the global gap}

\paragraph{Theorem O9 (compatible-idempotent certificate).}
Let $I$ be a finite index set and $(P_i)_{i\in I}$ pairwise commuting
idempotents. Set $M_i=\operatorname{ran}P_i$, $\Pi_i=\Pi_{M_i}$,
$Q_i=I-\Pi_i$, and $H=\sum_iQ_i$. Define
\[
 a_i=\|P_i-P_i^*\|,\quad
 \kappa_i=\sqrt{1+a_i^2},\quad
 L_i=\kappa_i+2\kappa_i^2a_i,
\]
and, for $i\ne j$,
\[
 b_{ij}=\min\left\{
 \frac{a_i}{\kappa_i},\frac{a_j}{\kappa_j},
 2\min(\kappa_i,\kappa_j)L_iL_j\|[P_i,P_j^*]\|
 \right\},\qquad b_{ii}=0.
\]
The entries are nonnegative. O4 proves symmetry of the adjoint-commutator
norm, so $B=(b_{ij})$ is a real symmetric matrix.
O5 and O6 imply $c(M_i,M_j)\leq b_{ij}$; O8 therefore gives
\[
 \boxed{\ \forall i\ne j,\ \forall\psi,\quad
 \operatorname{Re}\langle Q_i\psi,Q_j\psi\rangle
 \geq-b_{ij}\|Q_i\psi\|\|Q_j\psi\|.\ }
\]
This is exactly the hypothesis called hcos in an abstract projection-sum
formalization, now proved from compatible idempotents.

For completeness, for fixed $\psi$ put $x_i=\|Q_i\psi\|$.
Since the $Q_i$ are orthogonal projections,
\[
 \langle\psi,H\psi\rangle=\sum_i x_i^2,\qquad
 \langle\psi,H^2\psi\rangle
 =\sum_i x_i^2+2\sum_{i<j}
 \operatorname{Re}\langle Q_i\psi,Q_j\psi\rangle
 \geq x^{\mathsf T}(I-B)x.
\]
Diagonalizing the real symmetric matrix proves
$x^{\mathsf T}Bx\leq\rho(B)\|x\|_2^2$, since each eigenvalue is at most
its absolute value and hence at most the spectral radius.
It follows that $H^2\geq(1-\rho(B))H$.
Also $\ker H=\mathcal K:=\bigcap_iM_i$, because its quadratic form is
the sum of the nonnegative numbers $\|Q_i\psi\|^2$.
If $\alpha=1-\rho(B)>0$, every eigenvalue of the positive map $H$
satisfies $\lambda^2\geq\alpha\lambda$; hence it is zero or at least
$\alpha$. The spectral theorem therefore yields
\[
 H\geq(1-\rho(B))(I-\Pi_{\mathcal K}).
\]
This operator inequality remains meaningful when there is no excited
sector. For the empty index set, take $H=0$, $\mathcal K=\mathcal H$,
and the empty spectral radius to be zero.

\paragraph{Lemma O10 (weighted finite certificates).}
One may replace $B$ by any nonnegative symmetric matrix $\widehat B$
dominating the derived cosines. If $w_i>0$ and
$\sum_j\widehat b_{ij}w_j\leq rw_i$, then
$\rho(\widehat B)\leq r$.
Indeed $D_w^{-1}\widehat B D_w$ has nonnegative row sums at most $r$.
For an eigenvector choose a coordinate with maximal absolute value;
its eigenvalue equation and the triangle inequality give
$|\lambda|\leq r$. Similarity preserves eigenvalues.
The constant-vector choice gives the usual maximum-row-sum bound.
If $r<1$, the same quadratic-form proof gives gap at least $1-r$.

\subsection{Product filters and the explicit local neighborhood}

\paragraph{Lemma O11 (locality under similarity and orthogonalization).}
Let $\mathcal H=\bigotimes_{x\in V}\mathcal H_x$ with finite $V$ and
nonzero finite-dimensional site spaces. Suppose
$E_i=E_{i,S_i}\otimes I_{V\setminus S_i}$ are commuting orthogonal
projections. Given invertible site maps $L_x$, put
$L=\bigotimes_xL_x$ and $P_i=LE_iL^{-1}$.
Tensor multiplication outside $S_i$ cancels those factors, giving
\[
 P_i=(L_{S_i}E_{i,S_i}L_{S_i}^{-1})\otimes I_{V\setminus S_i}.
\]
Similarity preserves idempotence and commutativity.
The orthogonalizer O2 (or the tensor product of the range and its
orthogonal complement) gives
\[
 \Pi_{\operatorname{ran}P_i}
 =\Pi_{\operatorname{ran}(L_{S_i}E_{i,S_i}L_{S_i}^{-1})}
 \otimes I_{V\setminus S_i}.
\]
Thus the physical constraints have the same supports.
Disjoint supports imply $[P_i,P_j^*]=0$, hence $b_{ij}=0$.
Finally
\[
 \bigcap_i L\operatorname{ran}E_i
 =L\left(\bigcap_i\operatorname{ran}E_i\right).
\]
The reverse inclusion is immediate, and the forward inclusion follows
by applying $L^{-1}$ to a vector in every subspace. This identifies the
full, possibly degenerate, physical ground space.

\paragraph{Lemma O12 (local perturbation estimate).}
Suppose $L_x=e^{K_x}$ and $\|K_x\|\leq\varepsilon$ with $\varepsilon\geq0$.
Self-adjointness of $K_x$ is not needed for this bound.
For $|S_i|\leq s$ let $K_{S_i}$ be the sum of its embedded site generators.
These commute, so $L_{S_i}=e^{K_{S_i}}$ and
$\|K_{S_i}\|\leq s\varepsilon$.
For $E=E_{i,S_i}$ and $K=K_{S_i}$ put
$F(t)=e^{tK}Ee^{-tK}$ for $0\leq t\leq1$.
Differentiating the absolutely convergent matrix-exponential series gives
\[
 F'(t)=e^{tK}[K,E]e^{-tK}.
\]
The same series gives $\|e^{\pm tK}\|\leq e^{t\|K\|}$.
Since $\|E\|\leq1$, $\|[K,E]\|\leq2\|K\|$.
The fundamental theorem of calculus and the triangle inequality give
\[
 \|P_i-E_i\|
 \leq\int_0^1 2\|K\|e^{2t\|K\|}\,dt
 =e^{2\|K\|}-1
 \leq e^{2s\varepsilon}-1=:\delta.
\]
At $\|K\|=0$ both sides of the equality are zero directly.
Tensoring with the identity on a nonzero finite-dimensional factor
preserves the norm, justifying passage to the full system.
Since $E_i=E_i^*$,
\[
 a_i=\|(P_i-E_i)-(P_i-E_i)^*\|\leq2\delta.
\]
As $\kappa_i\geq1$, this implies $b_{ij}\leq2\delta$.

\paragraph{Corollary O13 (filtered-parent gap with derived hcos).}
Suppose every support has at most $s$ sites and intersects at most
$M$ other supports. O11 implies that each row of $B$ has at most $M$
nonzero off-diagonal entries; O12 bounds each by $2\delta$.
Consequently
$\rho(B)\leq2M(e^{2s\varepsilon}-1)$.
If this number is below one, O9 proves
\[
 H_L:=\sum_i(I-\Pi_{L\operatorname{ran}E_i})
 \ \geq\
 \left[1-2M(e^{2s\varepsilon}-1)\right]
 \left(I-\Pi_{\,L(\cap_i\operatorname{ran}E_i)}\right).
\]
In particular, whenever an excited sector exists, its spectral gap is
at least the bracketed number. This estimates local matrices and their
support incidences, never a condition number for the full product $L$.
Neither hcos nor support sparsity has been left as an unexplained input.

\paragraph{Optional nonuniform and adjoint-commutator bounds.}
Given individual bounds
$\|P_i-E_i\|\leq\delta_i$, the symmetric matrix
\[
 \widehat b_{ij}=
 \begin{cases}
 2\min(\delta_i,\delta_j),&i\ne j,\ S_i\cap S_j\ne\varnothing,\\
 0,&\text{otherwise}
 \end{cases}
\]
also satisfies hcos, by O5, O8, and O12. Apply O10 to it.
For a second optional estimate write $P_i=E_i+D_i$ with
$\|D_i\|\leq\delta_i$. Since the reference projections commute,
\[
 [P_i,P_j^*]=[E_i,D_j^*]+[D_i,E_j]+[D_i,D_j^*],
 \qquad
 \|[P_i,P_j^*]\|\leq2\delta_i+2\delta_j+2\delta_i\delta_j.
\]
This is valid, although generally weaker for the uniform neighborhood
than the separation term $a_i/\kappa_i$ already used above.
